\section{System Components}
\label{sec:system-components}

The Code Guardian prototype decomposes into a small set of components, each motivated by a specific design pressure derived from R1--R5. This section names them and states the rationale for each; the implementation detail (parameters, defensive-parsing logic, cache keys, retry policy) lives in Chapter~\ref{chap:implementation}, Section~\ref{sec:impl-components} and the appendices.

\paragraph{Context extraction (\ref{subsec:impl-context}).}
Real-time feedback in the IDE requires an analysis unit that is small enough to fit interactive latency budgets (R4) and large enough to support contextual reasoning (R1). Code Guardian therefore selects the smallest enclosing function as the default scope, with explicit selection, file, and workspace scopes for deeper inspection. Function scope keeps prompt size predictable; deeper program context (cross-file flows) is named as future work.

\paragraph{Knowledge retrieval (RAG) (\ref{subsec:impl-rag-knowledge}).}
A standalone LLM hallucinates vulnerability types and misclassifies severity \cite{ji2023hallucination}. Retrieval grounds reasoning in curated CWE and OWASP guidance \cite{lewis2020rag,karpukhin2020dpr} that evolves independently of model parameters. The corpus is local; only public metadata refresh (CVE/NVD summaries) ever leaves the device, satisfying R5 by separating code-bearing data paths from public-metadata update paths.

\paragraph{Vulnerability detection (\ref{subsec:impl-analyzer}).}
The detection component is the load-bearing element for R1 (accuracy) and R2 (consistency). It enforces a strict JSON output contract with format-level decoding, runs deterministic inference with multi-pass consensus filtering, and emits each finding with snippet-relative line indices. Detection and repair are split into two stages so each prompt focuses on a single task --- a separation that also preserves R3 by letting repair generation run only on findings the consensus filter accepts.

\paragraph{Repair generation and IDE integration.}
Repair suggestions are decision support, not autonomous edits. Each confirmed finding triggers a dedicated repair call returning a structured fix; the IDE integration layer surfaces the result as a Quick Fix that the developer can preview, apply, or discard --- preserving developer control (R3) and reversibility through the editor undo stack. IDE triggers (debounced edits, manual commands, workspace scans) are deliberately diverse so the system fits into both flow-state coding and explicit audit work (R4).

\paragraph{Supporting subsystems.}
Model management (runtime model switching), the analysis cache (LRU + TTL), the workspace dashboard (issue density and severity heuristics), and the vulnerability-data manager (cached public metadata) round out the component set. Each is independently testable and replaceable, which keeps the privacy boundary checkable in isolation. The next section explains how these components are orchestrated into workflows for real-time feedback, on-demand inspection, and batch scans.
